We are interested in cooperative aerial transport, in which $N$ multirotor uncrewed aerial vehicles (UAVs) lift a common payload via flexible cables~\cite{Michael2011, Fink2011, Bernard2011, Tagliabue2019}. Cable severance---from abrasion, hardware fatigue, or obstacle contact---is a structural fault: it discretely reduces the constraint set, redistributes load across surviving cables within one payload-pendulum period ($\tau_{\mathrm{pend}} = \SI{2.24}{\second}$ at $L = \SI{1.25}{\metre}$), and requires a compensatory response within that window.

Unilateral coupling, hybrid constraint-cardinality transitions at severance events~\cite{GoebelSanfeliceTeel2012, Brogliato2016}, and strict information locality together rule out classical cooperative and slung-load architectures: any admissible controller must deliver post-fault response without detection, annunciation, or reconfiguration and must remain analytically valid across structural transitions.

This paper shows that, for this plant class, a single feed-forward identity in the per-drone controller can provide the dominant passive fault-response mechanism, rather than relying on detection and reconfiguration: drone $i$'s altitude thrust feed-forward is set equal to its own measured rope tension,
\begin{equation}
  T_i^{\mathrm{ff}}(t) = T_i(t).
\end{equation}
While the use of measured interaction forces in control is not new, we show that, in this setting, this identity has a non-obvious consequence: it induces a certifiable stabilizing response across structural fault transitions on the plant's native timescale, without detection or reconfiguration.

When a peer's cable severs, the load redistributes within one payload-pendulum period ($\approx 2$~s) and the surviving drone's tension rises. Identity \eqref{eq:Tff} converts that rise at the next control tick into a proportional thrust increment, absorbing the new load at the plant's native rate without detection or communication. The mechanism sits inside a conventional cascade---an outer-loop PD controller with anti-swing slot shift, a single-step convex QP on the tilt-and-thrust envelope, and a geometric attitude loop---with \eqref{eq:Tff} as the sole nonstandard component.

We develop the consequences of \eqref{eq:Tff} along three axes. The primary contribution (C2, \cref{sec:stability}) is a hybrid practical-ISS theorem: under a single-active-regime QP doctrine and a spaced-fault dwell hypothesis $\tau_d \ge \tau_{\mathrm{pend}}$, the closed-loop tracking error admits a closed-form recovery envelope with Lyapunov decay rate $\lambda$ and per-fault-cycle contraction $\rho \in (0,1)$, conditional on bounded slack excursions, spaced fault events, and a single-active-regime QP---all audited in the simulation campaign. The strict $\rho < 1$ is achieved by \eqref{eq:Tff}: the dwell hypothesis is the minimum time over which $T_i^{\mathrm{ff}} = T_i$ restores the Lyapunov function below its pre-fault level (\cref{lem:dwell-decay}). The first secondary contribution (C1, \cref{sec:reduction-and-reference}) is a taut-cable reduction theorem: on a slack-excursion-bounded domain parameterized by the timescale ratio $\delta$ and duty-cycle bound $\eta$, an explicit first-order error budget justifies replacing the distributed rope dynamics with a lumped effective stiffness; all admissibility gates are met with margin by all six missions (\cref{tab:domain-audit}). The third contribution (C3, \cref{sec:extensions}) is a scalar admissible-window design rule for the $L_1$-adaptive altitude augmentation: two necessary conditions --- discretization stability and bandwidth coverage --- derived from the same scalar discrete pole together place the canonical gain $\Gamma$ inside a certifiable interval $[\Gamma_{\min}, \Gamma^\star]$, reducing reliance on empirical tuning. This is an auxiliary engineering rule, not a full discrete-time $L_1$ stability certificate for the complete interconnection.

The proposed controller is validated in high-fidelity Drake multibody simulation through a pre-declared six-variant capability demonstration and four targeted stress campaigns on a five-drone, \SI{10}{\kilo\gram}-payload, \SI{4}{\metre\per\second} Dryden-wind scenario under single and compound cable-severance faults. The feed-forward ablation campaign isolates identity \eqref{eq:Tff} as the causal mechanism: disabling it materially increases cruise-tracking error and post-fault altitude sag (\cref{tab:ff-ablation}). Two null findings are reported alongside: in these regimes, a composable extension layer is structurally inactive because the baseline already satisfies the relevant constraint. The architecture admits a sequence of $F \le N - 2$ unannounced severances under the actuator-margin envelope (\cref{sec:fault-and-controller}); the demonstrated $N=5$, $m_L=\SI{10}{\kilo\gram}$ admits $F \le 2$.

This paper identifies a stability mechanism for structural transitions that self-announces through actuator-side measurements. Four architectural consequences follow: \emph{locality} (no peer state, fault flag, or shared scheduler); \emph{passivity} (no detector, classifier, or reassignment rule); \emph{robustness to $F \le N - 2$ severances} by induction under the actuator-margin envelope; and \emph{team-size independence} (the per-drone law carries no $N$ structurally; dependence concentrates in the actuator-margin envelope). These properties are audited in \cref{sec:discussion-advance}; the two null findings (\cref{rem:mpc-null,rem:reshape-null}) confirm that the baseline mechanism explains the observed behavior.

\subsection{Related Work}
\label{sec:related-work}

Related work falls into four categories: centralized and distributed cooperative transport, geometric tracking and slung-load control, $L_1$ adaptive control, and fault-tolerant control.

\textit{Centralized and distributed cooperative transport} is a force-allocation architecture~\cite{Michael2011,Fink2011,Tagliabue2019} and distributed predictive variants~\cite{WangOng2017,LiLoianno2023,PallarLi2025,SuBhowmick2023} share a synchronous global-state requirement: the per-tick communication dependency is precisely what a cable-severance transition breaks at the actuator-tick rate, and none of these architectures accounts for cable severance in its stability argument. \textit{Geometric tracking}~\cite{LeeLeokMcClamroch2010}, \textit{flatness-based slung-load}~\cite{Sreenath2013}, and $\mathrm{SE}(3)$ \textit{payload control}~\cite{Lee2018} assume a fixed constraint set; multi-vehicle extensions~\cite{LeeSreenathKumar2013,SharmaSundaram2023} reintroduce global-state exchange, collapsing the strict locality the present architecture requires. \textit{$L_1$ adaptive control} provides matched-disturbance rejection with transient bounds~\cite{CaoHovakimyan2008,Hovakimyan2010book}; its extension to cooperative cable-fault settings under strict locality remains open~\cite{WuCheng2025}.

The fault-tolerant control literature is divided into two main categories: active architectures, which reconfigure themselves after diagnosing a fault~\cite{Blanke2015,ZhangJiang2008}, and passive architectures~\cite{Patton1997,Blanke2015}, which maintain robustness to absorb faults without reconfiguration\cite{LiuSuspensionFailure2021}. In the context of cooperative transport, approaches such as observer-based disturbance estimation with dynamic load reassignment address cable failures by detecting them and reconfiguring the system. Other methods, such as passive force control and leader-follower force estimation~\cite{NoCommTransport2021}, avoid the need to exchange information between UAVs but rely on a fixed number of cables.

This paper focuses specifically on the processing of structural transition announcements. In our approach, the measurement of tension on the actuator side acts as this announcement. This simplifies the system by eliminating the need for a lengthy process involving detection, classification, and reconfiguration. The key innovation presented here is the stability certification of a decentralized mechanism that converts tension into thrust in the presence of a structural fault. We directly measure the fault's impact via local cable tension.